% Part: model-theory
% Chapter: models-of-arithmetic
% Section: introduction

\documentclass[../../../include/open-logic-section]{subfiles}

\begin{document}

\olfileid{mod}{mar}{int}
\section{تعارف}

حساب دا \emph{معیاری نمونہ} اوہ
!!{structure}~$\Struct{N}$ اے جیہدے وچ $\Domain{N} = \Nat$
تے $\Obj{0}$، $\prime$، $+$، $\times$ تے $<$ دیاں تعبیراں
اُویں نیں جیویں توقع ہوندی اے۔ یعنی $\Obj{0}$ دی تعبیر $0$ اے،
$\prime$ جانشین فنکشن اے، $+$ قدرتی اعداد دی جمع تے $\times$
اوہناں دی ضرب اے، جتھے اعداد~$\Nat$ وچوں نیں۔ خاص طور اُتے،
\begin{align*}
  \Assign{\Obj{0}}{N} & = 0\\
  \Assign{\prime}{N}(n) & = n + 1\\
  \Assign{+}{N}(n, m) & = n + m\\
  \Assign{\times}{N}(n, m) & = nm
\end{align*}
بے شک $\Lang{L_A}$ دیاں ایسی ساختاں وی نیں جنہاں دے
دائرے~$\Nat$ توں وکھرے نیں۔ مثال لئی، اسیں
$\Struct{M}$ لے سکدے آں جیہدا دائرہ
$\Domain{M} = \{a\}^*$ اے۔
ایہ صرف اک علامت~$a$ توں بنیاں محدود ترتیب وار لڑیواں
دا سیٹ اے، یعنی $\emptyset$، $a$، $aa$، $aaa$، \dots؛
تعبیراں ایہ نیں:
\begin{align*}
  \Assign{\Obj{0}}{M} & = \emptyset\\
  \Assign{\prime}{M}(s) & = s \concat a\\
  \Assign{+}{M}(a^n, a^m) & = a^{n + m}\\
  \Assign{\times}{M}(a^n, a^m) & = a^{nm}
\end{align*}
چھوٹے ہون دا تعلق لڑیواں دیاں لمبائیاں دی قدرتی ترتیب دے مطابق اے۔
ایہ دو ساختاں اک خاص معنی وچ ”بنیادی طور اُتے اکّو جیہیاں“
نیں: فرق صرف اوہناں دے !!{domain}s دے !!{element}s وچ اے،
ایس گل وچ نہیں کہ تعبیر کرن والے فنکشن اوہناں !!{domain}s دے
!!{element}s نوں آپس وچ کِویں جوڑدے نیں۔ اسیں آکھدے آں کہ
ایہ دو !!{structure}s \emph{ہم شکل} نیں۔

محدود گواہی دے قضیے توں آسانی نال نکلدا اے کہ
$\Struct{N}$ وچ سچ کوئی وی نظریہ ایسے نمونے وی رکھدا اے
جو~$\Struct{N}$ نال ہم شکل نہیں۔ ایسی ساختاں نوں
\emph{غیر معیاری} آکھدے نیں۔ اوہناں دی دلچسپ گل ایہ اے کہ
معیاری نمونے دے !!{element}s، یعنی $\Struct{N}$ دے
یا اوہدے نال ہم شکل ساریاں !!{structure}s دے عنصر،
معیاری عددی عبارتاں~$\num{n}$ دیاں قدراں نال پورے ہو جاندے نیں؛
یعنی
\[
\Domain{N} = \Setabs{\Value{\num{n}}{N}}{n \in \Nat}
\]
پر رابنسن حساب (Q) دے غیر معیاری نمونے وچ ایہ گل درست نہیں:
جے $\Struct{M}$ غیر معیاری اے تے رابنسن حساب دے مسلّمے پورے کردا اے،
تاں گھٹ توں گھٹ اک
$x \in \Domain{M}$ ایسا اے کہ
$x \neq \Value{\num{n}}{M}$ ہر~$n$ لئی۔

ایہ غیر معیاری عنصر بڑے دلچسپ نیں: ایہ ”لامتناہی
قدرتی اعداد“ نیں۔ پر اک معنی وچ اوہناں دا وجود
نامکملیت دے مظاہر دی وجہ وی سمجھاؤندا اے۔ مثال لئی
پیانو حساب دی بے تضادی دا بیان~$\OCon[\Th{PA}]$
ویکھو، یعنی $\lnot
\lexists[x][\OPrf[\Th{PA}](x, \gn{\lfalse})]$۔
کیوں جے $\Th{PA}$ نہ $\OCon[\Th{PA}]$ نوں ثابت کردا اے
نہ $\lnot \OCon[\Th{PA}]$ نوں، ایس لئی دوواں وچوں
جیہڑا مرضی بیان $\Th{PA}$ وچ بے تضادی برقرار رکھ کے
شامل کیتا جا سکدا اے۔ کیوں جے $\Th{PA}$ بے تضاد اے،
$\Sat{N}{\OCon[\Th{PA}]}$، تے نتیجتاً
$\Sat/{N}{\lnot
  \OCon[\Th{PA}]}$۔ ایس لئی $\Struct{N}$،
$\Th{PA} \cup \{\lnot \OCon[\Th{PA}]\}$ دا
\emph{نمونہ نہیں}؛ ایس نظریے دے سارے نمونے لازماً
غیر معیاری نیں۔ $\Th{PA} \cup \{\lnot \OCon[\Th{PA}]\}$
دے نمونیاں وچ کوئی ایسا !!{element} ہونا چاہیدا اے
جو $\lexists[x][\OPrf[\Th{PA}](x, \gn{\lfalse})]$
نوں سچ بناؤن والا گواہ ہووے، یعنی نمونے دے اندر~$\Th{PA}$ توں
تضاد دی !!a{derivation} دا گوڈل عدد۔
ایہ !!a{element} معیاری نہیں ہو سکدا، کیوں جے
$\Th{PA} \Proves \lnot
\OPrf[\Th{PA}](\num{n}, \gn{\lfalse})$
ہر~$n$ لئی ہوندا اے۔

\end{document}
